% !TEX root = ../discretemath.tex

\section{Формула Кэли: доказательство через многопеременную производящую функцию}


\subsection{Многочлен остовных деревьев}

Рассмотрим полный граф $K_n$ на $\{1,\dots,n\}$.

\begin{Definition}{Многочлен остовных деревьев}{}
	\[
	F_n(x_1,\dots,x_n)=\sum_{T}\prod_{i=1}^{n}x_i^{\deg_i(T)-1},
	\]
	сумма по всем остовным деревьям $T$ графа $K_n$, $\deg_i(T)$ — степень вершины $i$ в $T$.
\end{Definition}

\begin{Remark}{}{}
	Все мономы одной степени:
	\[
	\sum_{i=1}^n(\deg_i(T)-1)=\sum_i\deg_i(T)-n=2(n-1)-n=n-2,
	\]
	так что $F_n$ — однородный многочлен степени $n-2$.
\end{Remark}

\begin{Example}{}{}
	При $n=6$ для дерева ниже степени $\deg_1=1,\deg_2=2,\deg_3=3,\deg_4=1,\deg_5=1,\deg_6=2$, моном $x_2^1x_3^2x_6^1$ степени $4=n-2$.
	\begin{center}
		\begin{tikzpicture}[v/.style={circle, draw, inner sep=2pt, minimum size=18pt, font=\small}]
			\node[v] (1) at (0, 1)    {$1$};
			\node[v] (3) at (2, 0)    {$3$};
			\node[v] (4) at (2,-1.4)  {$4$};
			\node[v] (2) at (4, 1)    {$2$};
			\node[v] (6) at (6, 1)    {$6$};
			\node[v] (5) at (8, 1)    {$5$};
			\draw (1)--(3)--(4)  (3)--(2)--(6)--(5);
		\end{tikzpicture}
	\end{center}
\end{Example}

\begin{Theorem}{Обобщённая формула Кэли}{}
	\[
	F_n(x_1,\dots,x_n)=(x_1+x_2+\cdots+x_n)^{n-2}.
	\]
\end{Theorem}

\begin{Corollary}{}{}
	При $x_1=\dots=x_n=1$: число остовных деревьев $K_n$ равно $n^{n-2}$ (классическая формула Кэли).
\end{Corollary}

\begin{proof}
	Индукция по $n$.

	\textbf{База ($n=2$):} единственное дерево — ребро $\{1,2\}$, $\deg_1=\deg_2=1$, $F_2=x_1^0x_2^0=1=(x_1+x_2)^0$.

	\textbf{Переход ($n-1\to n$):} по ИП $F_{n-1}=(x_1+\dots+x_{n-1})^{n-3}$.

	\emph{Шаг 1.} Для любого $k$ покажем $F_n(\dots,x_k{=}0,\dots)=\bigl(\sum_{i\neq k}x_i\bigr)^{n-2}$. Возьмём $k=n$. Множитель $0^{\deg_n(T)-1}$ ненулевой лишь при $\deg_n(T)=1$, то есть когда $n$ — лист, смежный с некоторой $j$. Тогда $\tilde T=T\setminus\{n\}$ — остовное дерево $K_{n-1}$, $\deg_j(T)=\deg_j(\tilde T)+1$, остальные степени не меняются, поэтому $\prod_i x_i^{\deg_i(T)-1}=x_j\prod_i x_i^{\deg_i(\tilde T)-1}$. Суммируя по парам $(\tilde T,j)$:
	\[
	F_n(x_1,\dots,x_{n-1},0)=\Bigl(\sum_{j=1}^{n-1}x_j\Bigr)F_{n-1}(x_1,\dots,x_{n-1})=(x_1+\dots+x_{n-1})^{n-2}.
	\]

	\emph{Шаг 2.} Оба многочлена однородны степени $n-2$. Для монома $m=x_1^{d_1}\cdots x_n^{d_n}$ с $\sum d_i=n-2<n$ найдётся $k$ с $d_k=0$. Тогда
	\[
	[m]\,F_n=[m]\,F_n(\dots,x_k{=}0,\dots)=[m]\Bigl(\sum_{i\neq k}x_i\Bigr)^{n-2}=[m]\,(x_1+\dots+x_n)^{n-2},
	\]
	где последнее равенство использует $d_k=0$. Коэффициенты при всех мономах совпадают, значит, многочлены равны.
\end{proof}

\begin{Remark}{}{}
	Доказательство — чисто алгебраическое сравнение коэффициентов многопеременной производящей функции $F_n$ с биномиальным (мультиномиальным) разложением $(x_1+\dots+x_n)^{n-2}$: коэффициент при $x_1^{d_1}\cdots x_n^{d_n}$ в правой части равен $\binom{n-2}{d_1,\dots,d_n}$ — это и есть число деревьев с заданными степенями вершин $\deg_i=d_i+1$ (формула Прюфера в скрытом виде).
\end{Remark}
